\documentclass[12pt,a4paper]{scrartcl}
\usepackage[utf8]{inputenc}
\usepackage[english,russian]{babel}
\usepackage{indentfirst}
\usepackage{graphicx}
\usepackage{amsmath}
\usepackage{amssymb}
\usepackage{listings}
\usepackage{xcolor}
\usepackage{geometry}

\geometry{left=2.5cm, right=1.5cm, top=2cm, bottom=2cm}

\lstset{
    language=Python,
    basicstyle=\ttfamily\small,
    keywordstyle=\color{blue}\bfseries,
    commentstyle=\color{gray}\itshape,
    stringstyle=\color{teal},
    showstringspaces=false,
    breaklines=true,
    frame=single,
    numbers=left,
    numberstyle=\tiny\color{gray},
    numbersep=5pt,
    tabsize=4
}

\begin{document}

\begin{titlepage}
    \begin{center}
        \large
        МИНИСТЕРСТВО НАУКИ И ВЫСШЕГО ОБРАЗОВАНИЯ РОССИЙСКОЙ ФЕДЕРАЦИИ

        Федеральное государственное бюджетное образовательное учреждение
        высшего образования

        \textbf{АДЫГЕЙСКИЙ ГОСУДАРСТВЕННЫЙ УНИВЕРСИТЕТ}
        \vspace{0.25cm}

        НОК <<Институт точных наук и цифровых технологий>>

        Кафедра управления в технических системах
        \vfill

        \textsc{Отчет по практике}\\[5mm]

        {\LARGE Программная реализация численного метода\\[3mm]
        \textit{Вычисление обратной матрицы методом Гаусса}}
        \bigskip

        1 курс, группа 1УТС
    \end{center}
    \vfill

    \newlength{\ML}
    \settowidth{\ML}{«\underline{\hspace{0.7cm}}» \underline{\hspace{2cm}}}
    \hfill\begin{minipage}{0.5\textwidth}
        Выполнил:\\
        \underline{\hspace{\ML}} В.\,Е.~Чешенко\\
        «\underline{\hspace{0.7cm}}» \underline{\hspace{2cm}} 2026 г.
    \end{minipage}%
    \bigskip

    \hfill\begin{minipage}{0.5\textwidth}
        Руководитель:\\
        \underline{\hspace{\ML}} С.\,В.~Теплоухов\\
        «\underline{\hspace{0.7cm}}» \underline{\hspace{2cm}} 2026 г.
    \end{minipage}%
    \vfill

    \begin{center}
        Майкоп, 2026 г.
    \end{center}
\end{titlepage}

\tableofcontents
\newpage

\section{Введение}
\label{sec:intro}

\subsection{Постановка задачи}

Для квадратной матрицы $A$ порядка $n$ требуется найти обратную матрицу $A^{-1}$
такую, что:
\begin{equation}
    A \cdot A^{-1} = A^{-1} \cdot A = E,
    \label{eq:def}
\end{equation}
где $E$~--- единичная матрица порядка~$n$.

Обратная матрица существует тогда и только тогда, когда матрица $A$ является
\textbf{невырожденной}, то есть $\det(A) \neq 0$.

\subsection{Теоретическая часть}

\subsubsection{Метод нахождения обратной матрицы}

Используется метод присоединения единичной матрицы (метод Гаусса-Жордана).
Идея основана на следующем наблюдении: если умножить обе части тождества
$A \cdot A^{-1} = E$ на $A^{-1}$ слева, то:
\[
    A^{-1} \cdot [A \,|\, E] = [E \,|\, A^{-1}].
\]

Таким образом, если к матрице $A$ справа приписать единичную матрицу того же
порядка и привести левую часть к единичной с помощью элементарных преобразований
строк, то правая часть автоматически станет $A^{-1}$.

\subsubsection{Алгоритм}

\begin{enumerate}
    \item \textbf{Построение сдвоенной матрицы.}
          К матрице $A$ размера $n \times n$ приписывается единичная матрица $E$:
          \[
              [A\,|\,E] =
              \begin{pmatrix}
                  a_{11} & \cdots & a_{1n} & 1 & \cdots & 0 \\
                  \vdots & \ddots & \vdots & \vdots & \ddots & \vdots \\
                  a_{n1} & \cdots & a_{nn} & 0 & \cdots & 1
              \end{pmatrix}.
          \]

    \item \textbf{Прямой и обратный ход Гаусса-Жордана.}
          Для каждого столбца $k = 1, \ldots, n$:
          \begin{enumerate}
              \item Выбрать ненулевой ведущий элемент в столбце $k$
                    (при необходимости переставить строки).
                    Если ненулевого нет~--- матрица вырождена, $A^{-1}$ не существует.
              \item Нормировать строку $k$: разделить все её элементы на
                    ведущий $a_{kk}$.
              \item Исключить $k$-й столбец из \emph{всех} остальных строк
                    (не только нижних):
                    \[
                        a_{ij} \leftarrow a_{ij} - a_{ik} \cdot a_{kj},
                        \quad i \neq k.
                    \]
          \end{enumerate}

    \item \textbf{Извлечение результата.}
          После $n$ итераций левая половина принимает вид $E$,
          а правая половина содержит $A^{-1}$:
          \[
              [A\,|\,E] \;\xrightarrow{\text{преобразования}}\; [E\,|\,A^{-1}].
          \]
\end{enumerate}

\subsubsection{Пример вручную}

Найдём обратную матрицу для:
\[
    A =
    \begin{pmatrix}
        2 & 1 \\
        5 & 3
    \end{pmatrix}.
\]

Строим сдвоенную матрицу:
\[
    \left(\begin{array}{cc|cc}
        2 & 1 & 1 & 0 \\
        5 & 3 & 0 & 1
    \end{array}\right).
\]

\textbf{Шаг 1} ($k=1$): нормируем строку~1 на $a_{11}=2$:
\[
    \left(\begin{array}{cc|cc}
        1 & 0{,}5 & 0{,}5 & 0 \\
        5 & 3     & 0     & 1
    \end{array}\right).
\]
Исключаем из строки~2 ($\text{строка}_2 - 5 \cdot \text{строка}_1$):
\[
    \left(\begin{array}{cc|cc}
        1 & 0{,}5 & 0{,}5  & 0 \\
        0 & 0{,}5 & -2{,}5 & 1
    \end{array}\right).
\]

\textbf{Шаг 2} ($k=2$): нормируем строку~2 на $a_{22}=0{,}5$:
\[
    \left(\begin{array}{cc|cc}
        1 & 0{,}5 & 0{,}5 & 0 \\
        0 & 1     & -5    & 2
    \end{array}\right).
\]
Исключаем из строки~1 ($\text{строка}_1 - 0{,}5 \cdot \text{строка}_2$):
\[
    \left(\begin{array}{cc|cc}
        1 & 0 &  3 & -1 \\
        0 & 1 & -5 &  2
    \end{array}\right).
\]

Результат:
\[
    A^{-1} =
    \begin{pmatrix}
         3 & -1 \\
        -5 &  2
    \end{pmatrix}.
\]

Проверка: $A \cdot A^{-1} = \begin{pmatrix}2&1\\5&3\end{pmatrix}
\begin{pmatrix}3&-1\\-5&2\end{pmatrix} =
\begin{pmatrix}1&0\\0&1\end{pmatrix} = E$. \checkmark

\section{Ход работы}
\label{sec:work}

\subsection{Описание программы}

Программа реализована на языке Python. Основные функции:
\begin{itemize}
    \item \texttt{input\_int(prompt, min\_val, max\_val)} ---
          ввод целого числа с проверкой диапазона;
    \item \texttt{input\_float(prompt)} ---
          ввод вещественного числа с обработкой ошибок;
    \item \texttt{input\_matrix(n)} ---
          построчный ввод квадратной матрицы $n \times n$;
    \item \texttt{print\_matrix(matrix, title)} ---
          форматированный вывод матрицы;
    \item \texttt{inverse\_matrix(matrix)} ---
          реализация алгоритма; возвращает $A^{-1}$ или
          \texttt{None} при вырожденной матрице;
    \item \texttt{mat\_mul(a, b)} ---
          перемножение двух матриц для проверки;
    \item \texttt{verify(original, inv)} ---
          проверка $A \cdot A^{-1} = E$ с выводом максимального отклонения.
\end{itemize}

\subsection{Контроль вводимых данных}

\begin{enumerate}
    \item Порядок матрицы $n$ должен быть \emph{целым числом} от~1 до~8.
          При ошибке выдаётся сообщение и запрашивается повторный ввод.
    \item Элементы матрицы должны быть \emph{вещественными числами}.
          Нечисловой ввод обрабатывается без аварийного завершения.
    \item Если в процессе преобразований ведущий элемент оказывается нулём
          во всём столбце, программа сообщает о вырожденности матрицы
          и завершает работу корректно.
\end{enumerate}

\subsection{Код программы}
\label{sec:code}

\begin{lstlisting}[caption={inverse\_matrix.py --- ключевые функции}]
import copy

def inverse_matrix(matrix):
    n = len(matrix)
    # Build augmented matrix [A | E]
    aug = []
    for i in range(n):
        row = copy.copy(matrix[i])
        for j in range(n):
            row.append(1.0 if i == j else 0.0)
        aug.append(row)

    for col in range(n):
        # Find non-zero pivot
        pivot_row = None
        for row in range(col, n):
            if abs(aug[row][col]) > 1e-12:
                pivot_row = row
                break
        if pivot_row is None:
            return None  # singular matrix

        # Swap rows if needed
        if pivot_row != col:
            aug[col], aug[pivot_row] = aug[pivot_row], aug[col]

        # Normalize pivot row
        pivot = aug[col][col]
        aug[col] = [x / pivot for x in aug[col]]

        # Eliminate from ALL other rows
        for row in range(n):
            if row == col:
                continue
            factor = aug[row][col]
            aug[row] = [aug[row][k] - factor * aug[col][k]
                        for k in range(2 * n)]

    # Extract right half -> A^(-1)
    return [aug[i][n:] for i in range(n)]


def mat_mul(a, b):
    n = len(a)
    return [[sum(a[i][k] * b[k][j] for k in range(n))
             for j in range(n)] for i in range(n)]
\end{lstlisting}

\subsection{Пример работы программы}

Найдём обратную матрицу для:
\begin{equation}
    A =
    \begin{pmatrix}
        1 & 2 & 3 \\
        0 & 1 & 4 \\
        5 & 6 & 0
    \end{pmatrix}.
    \label{eq:example}
\end{equation}

\begin{table}[h]
    \centering
    \caption{Исходная матрица~\eqref{eq:example}}
    \label{tab:input}
    \begin{tabular}{|c|c|c|}
        \hline
        1 & 2 & 3 \\
        \hline
        0 & 1 & 4 \\
        \hline
        5 & 6 & 0 \\
        \hline
    \end{tabular}
\end{table}

Вывод программы:

\begin{verbatim}
Обратная матрица A^(-1):
+------------+------------+------------+
| -24.00000  |  18.00000  |   5.00000  |
|  20.00000  | -15.00000  |  -4.00000  |
|  -5.00000  |   4.00000  |   1.00000  |
+------------+------------+------------+

Проверка A * A^(-1) = E:
  Произведение A * A^(-1):
+------------+------------+------------+
|  1.00000   |  0.00000   |  0.00000   |
|  0.00000   |  1.00000   |  0.00000   |
|  0.00000   |  0.00000   |  1.00000   |
+------------+------------+------------+

  Максимальное отклонение от E: 0.00e+00
  Результат верен.
\end{verbatim}

\subsection{Обработка особых случаев}

При вводе вырожденной матрицы, например:
\[
    B =
    \begin{pmatrix}
        1 & 2 & 3 \\
        4 & 5 & 6 \\
        7 & 8 & 9
    \end{pmatrix}
    \quad (\det B = 0)
\]
программа выводит:
\begin{verbatim}
Обратная матрица не существует
(матрица вырождена: определитель равен нулю).
\end{verbatim}

\section{Заключение}

В ходе практической работы реализована программа на языке Python,
вычисляющая обратную матрицу методом Гаусса-Жордана. В программе обеспечены:

\begin{itemize}
    \item построение сдвоенной матрицы $[A\,|\,E]$ и приведение её к виду $[E\,|\,A^{-1}]$;
    \item обнаружение вырожденных матриц с корректным завершением работы;
    \item проверка результата через произведение $A \cdot A^{-1}$
          с выводом максимального отклонения от~$E$;
    \item контроль корректности всех вводимых данных.
\end{itemize}

Метод протестирован на матрицах различных порядков, включая вырожденный и
тривиальный случаи. Результаты совпали с эталонными значениями с машинной точностью.

\begin{thebibliography}{9}
    \bibitem{samarsky}
        Самарский А.\,А., Гулин А.\,В.
        \textit{Численные методы.} ---
        М.: Наука, 1989. --- 432~с.

    \bibitem{bahvalov}
        Бахвалов Н.\,С., Жидков Н.\,П., Кобельков Г.\,М.
        \textit{Численные методы.} ---
        М.: БИНОМ, 2008. --- 636~с.

    \bibitem{python}
        Van Rossum G., Drake F.\,L.
        \textit{Python 3 Reference Manual.} ---
        Scotts Valley, CA: CreateSpace, 2009.
\end{thebibliography}

\end{document}
